\documentclass[10pt,a4paper]{amsart}
\usepackage[margin=19mm]{geometry}
\usepackage{amsmath,amssymb,mathtools}
\usepackage[scr=rsfs,cal=euler]{mathalfa}
\usepackage{xfrac}
\usepackage[unicode,hidelinks,bookmarks=false]{hyperref}
\newcommand{\ie}{i.e.\ }
\newcommand{\cf}{cf.\ }
\newcommand{\resp}{respectively\ }
\newcommand{\Subsection}{\subsection}
\providecommand{\nobreakdash}{\nobreak\hbox{-}}
\setlength{\emergencystretch}{2em}
\multlinegap0pt
\usepackage{bm}
\usepackage{bbm}
\usepackage{dsfont}
\usepackage{xspace}
\usepackage{cancel}
\usepackage{mathtools}
\usepackage{amsthm}
\makeatletter
\@ifundefined{Prop}{\newtheorem{Prop}{Proposition}}{}
\makeatother
\makeatletter
\expandafter\ifx\csname acknowledgements\endcsname\relax
\newenvironment{acknowledgements}{%
	\titlepage
	\null\vfil
	\@beginparpenalty\@lowpenalty
	\begin{center}%
		\bfseries \ackname
		\@endparpenalty\@M
\end{center}}%
\fi
\DeclarePairedDelimiter{\scal}{\langle}{\rangle}
\newcommand\uxi{{\underline \xi}}
\newcommand\vxi{{\vec \xi}}
\makeatother
\title[A Kac system interacting with two heat reservoirs: the shear]{A Kac system interacting with two heat reservoirs: the shearing case}
\author{Federico Bonetto, Matthew Powell}
\date{Source: arXiv:2606.30529v1 (CC BY 4.0)}
\begin{document}
\maketitle

\noindent\emph{Source and licence.} A Kac system interacting with two heat reservoirs: the shearing case. Version arXiv:2606.30529v1 (CC BY 4.0), distributed under the Creative Commons Attribution 4.0 International licence, as recorded in the arXiv licence metadata for this exact version and shown on the abstract page https://arxiv.org/abs/2606.30529v1. Full rights notice: licenses/arxiv-2606.30529-CC-BY-4.0.txt.\par\smallskip

\noindent\textit{Editorial scope.} Counted unit: the Prop whose source label is \texttt{Prop:tedious} in the article (source0.tex lines 1177--1184, source label \texttt{Prop:tedious}), delivered together with the article's own complete original proof of it (source0.tex lines 1186--1260, one \texttt{proof} environment of 75 lines). No mathematics is written, repaired or re-derived: statement and proof are the article's own text line for line. Required context retained uncounted: eq:IVPM1 (lines 389--408). Every definition, display and label of those lines is retained unchanged. Omitted from this package and not redistributed: 1--388, 409--1176, 1185--1185, 1261--1793. Nothing inside a retained excerpt is omitted or altered: every retained body is delivered exactly as the article's lines are written, so no omission receipt of any kind accompanies this package. Citations: the retained excerpts carry no citation command at all, so no bibliography entry of the article is transcribed and no reference is redistributed. Reference adaptation: none was needed and none was made; every reference inside the delivered unit resolves inside it, so no unresolved reference or citation is produced. Printed slips kept literal: no printed word, symbol, hypothesis or number of the counted statement or of its proof was altered. Layout and class: the article's own class is \texttt{article} with packages including amsmath, amsbsy, amsfonts, amssymb, amsthm, amscd, amsxtra, amstext, amsopn, mathtools, emptypage, dsfont, yfonts, latexsym; the wrapper is the lane's pinned \texttt{amsart} editorial wrapper at 10pt on A4, which restates the theorem-like environments the retained text needs and adds the article's own macro definitions that the retained text needs (\texttt{\textbackslash scal}, \texttt{\textbackslash uxi}, \texttt{\textbackslash vxi}), with extra wrapper packages (bm, bbm, dsfont, xspace, cancel, mathtools). The delivered numbering is the unit's own. Assets: no retained span contains a figure environment, a graphics-inclusion command, a PDF-page inclusion, a table or a TikZ picture; no image file is copied into this package, the package declares no asset dependency and no image file is redistributed with it. Licence: version arXiv:2606.30529v1 of this article is distributed under the Creative Commons Attribution 4.0 International licence, as recorded in the arXiv licence metadata for this exact version and shown on the abstract page https://arxiv.org/abs/2606.30529v1. A full rights notice is delivered at licenses/arxiv-2606.30529-CC-BY-4.0.txt. Characters: every retained excerpt is pure ASCII, so the delivered engine, XeLaTeX with its default Latin Modern text font, reports no missing character.
\begin{equation}\label{eq:IVPM1}
\frac{d}{dt}\begin{pmatrix}
		\vec m_+ \\ \vec m_S \\ \vec m_-
	\end{pmatrix}=
\frac 13\begin{pmatrix}
-\frac{M}{N} & \frac MN & 0\\
1 & -2 & 1\\
0 & \frac{M}{N} & -\frac MN
\end{pmatrix}
\begin{pmatrix}
		\vec m_+ \\ \vec m_S \\ \vec m_-
	\end{pmatrix},
\qquad\qquad
\begin{pmatrix}
		\vec m_+(0) \\ \vec m_S(0) \\ \vec m_-(0)
	\end{pmatrix}=
\begin{pmatrix}
		\vec p_+ \\ \vec p_S \\ \vec p_-
	\end{pmatrix}
\end{equation}

\begin{Prop}\label{Prop:tedious}
For all $s' \geq 0$ we have 
\begin{equation*} 
d_2\left(\Psi(s') \Gamma^M_{\infty,s'},
\frac{d}{ds'} \Gamma^M_{\infty,s'}\right) \leq \frac 1 6  \sum_{\sigma\in\{+,-\}}\big(|\tau_\sigma(s') - \tau_\infty| + 
2\|\vec\delta_\sigma(s')\|^2\big)
\end{equation*} 
\end{Prop}

\begin{proof}
We begin by observing that
\begin{align*}
\widehat\Psi(s') \widehat\Gamma^M_{\infty,s'}(\uxi) - &\widehat{\frac{d}{ds'}\Gamma^M_{\infty,s'}}(\uxi) =\\
& \sum_{k=1}^M \Gamma_{\infty,s'}^{M - 1}(\uxi^k) 
\bigg(\int_{\mathds S^2}\left(e^{i\langle\vec \delta_+(s'),\vec 0^*_k\rangle}\widehat \Gamma_{\tau_+(s')}(\vec
0^*_k)+e^{i\langle\vec \delta_-(s'),\vec 0^*_k\rangle}\widehat \Gamma_{\tau_-(s')}(\vec
0^*_k)\right)\widehat \Gamma_\infty(\vxi_k^*) -
\\ 
&\quad 2\widehat \Gamma_\infty(\vxi_k) - i\langle \dot{\vec m}_S(s'), \vxi_k\rangle \Gamma_\infty(\vxi_k)\bigg) =: 
\sum_{k=1}^M \widehat\Gamma_{\infty,s'}^{M - 1}(\uxi^k) D_2(\vxi_k)
\end{align*}
where $0^*_k=\langle\omega,\vxi_k\rangle\omega$.
It will suffices to understand $D_2(\vxi)$.
We begin by writing $\widehat \Gamma_{\tau_\sigma(s')}(0^*) = \widehat \Gamma_{\tau_\sigma(s')}(0^*) - \widehat 
\Gamma_{\infty}(0^*) + \widehat \Gamma_{\infty}(0^*)$ to obtain
\begin{align*}
|D_2(\vxi)|&= \bigg|\int_{\mathds S^2} e^{i\langle\vec \delta_+(s'),\vec 0^*\rangle}( \widehat 
\Gamma_{\tau_+(s')}(\vec 0^*)  - \widehat \Gamma_{\infty}(0^*))
\widehat\Gamma_{\infty}(\vxi^*) \, d\omega 
 + 
\int_{\mathds S^2} e^{i\langle\vec \delta_+(s'),\vec 0^*\rangle}\widehat \Gamma_{\infty}(\vxi) \, d\omega
\\
&\quad +
\int_{\mathds S^2} e^{i\langle\vec \delta_-(s'),\vec 0^*\rangle}( \widehat \Gamma_{\tau_-(s')}(\vec 0^*)  - \widehat 
\Gamma_{\infty}(0^*))
\widehat\Gamma_{\infty}(\vxi^*) \, d\omega 
 + 
\int_{\mathds S^2} e^{i\langle\vec \delta_-(s'),\vec 0^*\rangle}\widehat \Gamma_{\infty}(\vxi) \, d\omega
\\
&\quad -
2\widehat \Gamma_\infty(\vxi) - i\langle \dot{m}_S(s'), \vxi\rangle \widehat\Gamma_\infty(\vxi) \bigg| 
\\
&\leq 
\int_{\mathds S^2} \left| \widehat \Gamma_{\tau_+(s')}(\vec 0^*)  - \widehat \Gamma_{\infty}(0^*)\right|
 \, d\omega + \int_{\mathds S^2} \left| \widehat \Gamma_{\tau_-(s')}(\vec 0^*)  - \widehat \Gamma_{\infty}(0^*)\right|
 \, d\omega
 \\
 &\quad + 
 \left|\int_{\mathds S^2} e^{i\langle\vec \delta_+(s'),\vec 0^*\rangle} + e^{i\langle\vec \delta_-(s'),\vec 0^*\rangle} 
 \, d\omega - 2 - i\langle \dot{m}_S(s'), \vxi\rangle \right|
\end{align*}
Using that $\int_{\mathds S^2}\|0^*\|^2d\omega=\frac13\|\vxi\|^2$, the first term in this estimate satisfies
\[
 \int_{\mathds S^2} \left| \widehat \Gamma_{\tau_\sigma(s')}(\vec 0^*)  - \widehat \Gamma_{\infty}(0^*)\right|
 \, d\omega \leq \frac 1 3 d_2(\Gamma_{\tau_\sigma(s')}, \Gamma_\infty)\|\vxi\|^2 \leq \frac 1 6 |\tau_\sigma(s') - 
 \tau_\infty|\|\vxi\|^2 \, .
\]
We can control the second term by expanding $e^{i\scal{\vec \delta_\sigma(s'),\vec 0^*}}$ using 
\begin{equation*}
e^{x} = 1 + x + x^2\int_0^1 e^{tx}(1 - t) \ dt
\end{equation*}
and using \eqref{eq:IVPM1} to obtain
\begin{align*}
\bigg|\int_{\mathds
S^2} &e^{i\langle\vec \delta_+(s'),\vec 0^*\rangle} + e^{i\langle\vec \delta_-(s'),\vec 0^*\rangle}
-2 - i\langle \dot{m_S}, \xi\rangle \bigg|
\\
&=
\bigg|
\int_{\mathds S^2} \langle\vec \delta_+(s'),\vec 0^*\rangle^2 \int_0^1 e^{i\langle\vec \delta_+(s'),\vec 0^*\rangle t} 
(1 - 
t)\, dt \, d\omega + 
\int_{\mathds S^2} \langle\vec \delta_-(s'),\vec 0^*\rangle^2\int_0^1 e^{i\langle\vec \delta_-(s'),\vec 0^*\rangle t} 
(1 - t)\, 
dt \, d\omega \bigg|
\\
&\leq 
\int_{S^2} \langle\vec \delta_+(s'),\vec 0^*\rangle^2 + \langle\vec \delta_-(s'),\vec 0^*\rangle^2 \, d\omega 
\leq
\frac 1 3 \left(\|\vec \delta_+(s')\|^2  + \|\vec \delta_-(s')\|^2\right) \|\xi\|^2 
\end{align*}

Combining these estimates we finally get the thesis.
\end{proof}

\end{document}
